% !TEX spellcheck = en_US
% !TEX spellcheck = LaTeX
\documentclass[letterpaper,10pt,english]{article}
\usepackage{%
amsfonts,%
amsmath,%
amssymb,%
amsthm,%
babel,%
bbm,%
%biblatex,%
caption,%
centernot,%
color,%
enumerate,%
%enumitem,%
epsfig,%
epstopdf,%
etex,%
fancybox,%
framed,%
fullpage,%
%geometry,%
graphicx,%
hyperref,%
latexsym,%
mathptmx,%
mathtools,%
multicol,%
paralist,%
pgf,%
pgfplots,%
pgfplotstable,%
pgfpages,%
proof,%
psfrag,%
%subfigure,%
tcolorbox,%
tfrupee,%
tikz,%
times,%
%ulem,%
url,%
xcolor,%
mathpazo,%
}

\definecolor{shadecolor}{gray}{.95}%{rgb}{1,0,0}
\usepackage[margin=1in,top=0.75in]{geometry}
\usepackage[mathscr]{eucal}
\usepgflibrary{shapes}
\usepgfplotslibrary{fillbetween}
\usetikzlibrary{%
arrows,%
backgrounds,%
chains,%
decorations.pathmorphing,% /pgf/decoration/random steps | erste Graphik
decorations.text,% 
matrix,%
positioning,% wg. " of "
fit,%
patterns,%
petri,%
plotmarks,%
scopes,%
shadows,%
shapes.misc,% wg. rounded rectangle
shapes.arrows,%
shapes.callouts,%
shapes%
}

%\pgfplotsset{compat=newest} %<------ Here
\pgfplotsset{compat=1.11} %<------ Or use this one

\theoremstyle{plain}
\newtheorem{thm}{Theorem}[section]
\newtheorem{lem}[thm]{Lemma}
\newtheorem{prop}[thm]{Proposition}
\newtheorem{cor}[thm]{Corollary}
\newtheorem{clm}[thm]{Claim}

\theoremstyle{definition}
\newtheorem{defn}[thm]{Definition}
\newtheorem{conj}[thm]{Conjecture}
\newtheorem{exmp}[thm]{Example}
\newtheorem{axiom}[thm]{Axiom}
\newtheorem{assum}[thm]{Assumption}
\newtheorem{exerc}[thm]{Exercise}
\newtheorem*{defin}{Definition}

\theoremstyle{remark}
\newtheorem{rem}{Remark}
\newtheorem{note}{Note}

\newcommand{\iid}{\emph{i.i.d.}\ }
\newcommand{\Cov}{\operatorname{Cov}}
%\newcommand{\det}{\operatorname{det}}
%\newcommand{\Real}{\mathbb{R}}
\newcommand{\tr}{\operatorname{tr}}
\newcommand{\Var}{\operatorname{Var}}
\newcommand{\cov}{\operatorname{cov}}

%\renewcommand{\proof}[1]{\begin{proof}#1\end{proof}}
\newcommand{\EQ}[1]{\begin{equation*}#1\end{equation*}}
\newcommand{\EQN}[1]{\begin{equation}#1\end{equation}}
\newcommand{\eq}[1]{\begin{align*}#1\end{align*}}
\newcommand{\meq}[2]{\begin{xalignat*}{#1}#2\end{xalignat*}}
\newcommand{\norm}[1]{\left\lVert#1\right\rVert}
\newcommand{\inner}[1]{\left\langle #1\right\rangle}
\newcommand{\abs}[1]{\left\lvert#1\right\rvert}
\newcommand{\expect}[1]{\mathbb{E}\left[{#1}\right]}
\newcommand{\prob}[1]{\mathbb{P}\left[{#1}\right]}
\newcommand{\given}{\; \big\vert \;} 
\newcommand{\set}[1]{\left\{#1\right\}} 
\newcommand{\indicator}[1]{1_{\set{#1}}} 
\newcommand{\Ind}[1]{\mathbbm{1}_{#1}} 
\newcommand{\SetIn}[1]{\mathbbm{1}_{\set{#1}}} 
\newcommand{\red}[1]{\textcolor{red}{#1}} 
\newcommand{\floor}[1]{\left\lfloor#1\right\rfloor}
\newcommand{\ceil}[1]{\left\lceil#1\right\rceil}


\newcommand{\C}{\mathbb{C}}
\newcommand{\D}{\mathbb{D}}
\newcommand{\E}{\mathbb{E}}
\newcommand{\F}{\mathbb{F}}
\newcommand{\N}{\mathbb{N}}
\renewcommand{\P}{\mathbb{P}}
\newcommand{\Q}{\mathbb{Q}}
\newcommand{\R}{\mathbb{R}}
\newcommand{\Z}{\mathbb{Z}}

\newcommand{\bA}{\mathbf{A}}
\newcommand{\bI}{\mathbf{I}}
\newcommand{\bU}{\mathbf{1}}
\newcommand{\bx}{\mathbf{x}}
\newcommand{\bX}{\mathbf{X}}
\newcommand{\bZ}{\mathbf{Z}}


\newcommand{\cA}{\mathcal{A}}
\newcommand{\cB}{\mathcal{B}}
\newcommand{\cC}{\mathcal{C}}
\newcommand{\cF}{\mathcal{F}}
\newcommand{\cG}{\mathcal{G}}
\newcommand{\cM}{\mathcal{M}}
\newcommand{\cN}{\mathcal{N}}
\newcommand{\cP}{\mathcal{P}}
\newcommand{\cT}{\mathcal{T}}
\newcommand{\cX}{\mathcal{X}}
\newcommand{\cY}{\mathcal{Y}}

\newcommand{\sA}{\mathscr{A}}
\newcommand{\sB}{\mathscr{B}}
\newcommand{\sC}{\mathscr{C}}
\newcommand{\sE}{\mathscr{E}}
\newcommand{\sF}{\mathscr{F}}
\newcommand{\sG}{\mathscr{G}}
\newcommand{\sH}{\mathscr{H}}
\newcommand{\sL}{\mathscr{L}}
\newcommand{\sS}{\mathscr{S}}
\newcommand{\sT}{\mathscr{T}}
\newcommand{\sX}{\mathscr{X}}
\newcommand{\sY}{\mathscr{Y}}
\newcommand{\sZ}{\mathscr{Z}}

\newcommand{\tS}{\tilde{S}}
% Debug
\newcommand{\todo}[1]{\begin{color}{blue}{{\bf~[TODO:~#1]}}\end{color}}

% a few handy macros

\renewcommand{\le}{\leqslant}
\renewcommand{\ge}{\geqslant}
\newcommand\matlab{{\sc matlab}}
\newcommand{\goto}{\rightarrow}
\newcommand{\bigo}{{\mathcal O}}
%\newcommand{\half}{\frac{1}{2}}
%\newcommand\implies{\quad\Longrightarrow\quad}
\newcommand\reals{{{\rm l} \kern -.15em {\rm R} }}
\newcommand\complex{{\raisebox{.043ex}{\rule{0.07em}{1.56ex}} \hskip -.35em {\rm C}}}


% macros for matrices/vectors:

% matrix environment for vectors or matrices where elements are centered
\newenvironment{mat}{\left[\begin{array}{ccccccccccccccc}}{\end{array}\right]}
\newcommand\bcm{\begin{mat}}
\newcommand\ecm{\end{mat}}

% matrix environment for vectors or matrices where elements are right justifvied
\newenvironment{rmat}{\left[\begin{array}{rrrrrrrrrrrrr}}{\end{array}\right]}
\newcommand\brm{\begin{rmat}}
\newcommand\erm{\end{rmat}}

% for left brace and a set of choices
%\newenvironment{choices}{\left\{ \begin{array}{ll}}{\end{array}\right.}
\newcommand\when{&\text{if~}}
\newcommand\otherwise{&\text{otherwise}}
% sample usage:
%\delta_{ij} = \begin{choices} 1 \when i=j, \\ 0 \otherwise \end{choices}


% for labeling and referencing equations:
\newcommand{\eql}{\begin{equation}\label}
\newcommand{\eqn}[1]{(\ref{#1})}
% can then do
%\eql{eqnlabel}
%...
%\end{equation}
% and refer to it as equation \eqn{eqnlabel}.
% some useful macros for finite difference methods:
\newcommand\unp{U^{n+1}}
\newcommand\unm{U^{n-1}}
% for chemical reactions:
\newcommand{\react}[1]{\stackrel{K_{#1}}{\rightarrow}}
\newcommand{\reactb}[2]{\stackrel{K_{#1}}{~\stackrel{\rightleftharpoons}
 {\scriptstyle K_{#2}}}~}
\makeatletter
\def\th@plain{%
\thm@notefont{}% same as heading font
\itshape % body font
}
\def\th@definition{%
\thm@notefont{}% same as heading font
\normalfont % body font
}
\makeatother
\date{}


\title{Lecture-17: Tractable Random Processes}
\author{}

\begin{document}
\maketitle

\section{Examples of Tractable Stochastic Processes} 
Recall that a random process $X: \Omega \to \sX^T$ defined on the probability space $(\Omega,\sF, P)$ with index set $T$ and state space $\sX \subseteq \R$,  
is completely characterized by its finite dimensional distributions $F_{X_S}: \R^S \to [0,1]$ for all finite $S \subseteq T$, 
where 
\EQ{
F_{X_S}(x_S) 
\triangleq P(A_{X_S}(x_S))
= P(\cap_{s \in S}X_s^{-1}(-\infty, x_s]),\quad x_S \in \R^S. 
}
Simpler characterizations of a stochastic process $X$ are in terms of its moments. 
That is, the first moment such as mean, and the second moment such as correlations and covariance functions. 
\meq{3}{
&m_X(t) \triangleq \E X_t,&
&R_X(t,s) \triangleq \E X_tX_s,&
&C_X(t,s) \triangleq \E (X_t - m_X(t))(X_s-m_X(s)).
}
In general, it is very difficult to characterize a stochastic process completely in terms of its finite dimensional distribution. 
However, we have listed few analytically tractable examples below, where we can completely characterize the stochastic process. 

%We will consider the probability space $(\Omega, \sF, P)$, and a random process $X: \Omega \to \sX^T$ for index set $T$ and state space $\sX \subseteq \R$. 
\subsection{Independent and identically distributed (\iid) processes} 
\begin{defn}[\iid process]
A random process $X: \Omega \to \sX^T$ is an \textbf{independent and identically distributed} (\iid) random process with the common distribution $F:\R \to [0,1]$, if for any finite $S \subseteq T$ and a real vector $x_S \in \R^S$ we can write the finite dimensional distribution for this process as 
\EQ{
F_{X_S}(x_S) = P\left(\cap_{s \in S}\set{X_s(\omega) \le x_s}\right) = \prod_{s \in S}F(x_s).
}
\end{defn}
\begin{rem}
It's easy to verify that the first and the second moments are independent of time indices. 
That is, if $0 \in T$ then $X_t = X_0$ in distribution, and we have 
\meq{3}{
&m_X = \E X_0 , && R_X(t,s) = (\E X_0^2)\SetIn{t=s} + m_X^2\SetIn{t\neq s}, && C_X(t,s) = \Var(X_0)\SetIn{t=s}.
}
\end{rem}

\subsection{Stationary processes}
\begin{defn}[Stationary process]
We consider the index set $T \subseteq \R$ that is closed under addition and subtraction.   
A stochastic process $X: \Omega \to \sX^T$ is \textbf{stationary} if all finite dimensional distributions are shift invariant. 
That is, for any finite $S \subseteq T$ and $t \in T$, we have 
\EQ{
F_{X_S}(x_S) 
= P(\cap_{s\in S}\set{X_s(\omega) \le x_s}) 
= P(\cap_{s\in S}\set{X_{s+t}(\omega) \le x_s}) 
= F_{X_{t+S}}(x_S).
}
\end{defn}
\begin{rem} 
That is, for any finite $n \in \N$ and $t \in \R$, the random vectors $(X_{s_1}, \dots, X_{s_n})$ and $(X_{s_1+t}, \dots, X_{s_1+t})$ have the identical joint distribution for all $s_1 \le \dots \le s_n$.  
\end{rem}

\begin{shaded}
%\begin{exmp}[IID processes] 
\begin{lem} 
Any \emph{i.i.d.} process with index set $T \subseteq \R$ is stationary.
\end{lem}
\begin{proof}
Let $X: \Omega \to \sX^T$ be an \emph{i.i.d.} random process, where $T \subseteq \R$. 
Then, for any finite index subset $S \subseteq T, t \in T$ and $x_S \in \R^S$, we can write 
\EQ{
F_{X_S}(x_S) 
= P(\bigcap_{s \in S}\set{X_s \le x_s}) 
= \prod_{s \in S}P\circ X_s^{-1}(-\infty, x_s] 
= \prod_{s \in S}P\circ X_{t+s}^{-1}(-\infty, x_s] 
%&= \prod_{u \in t+S}P\set{X_u \le x_u} 
= P(\bigcap_{s \in t+S}\set{X_s \le x_s})  
= F_{X_{t+S}}(x_S).
}
First equality follows from the definition, the second from the independence of process $X$, 
the third from the identical distribution for the process $X$. 
In particular, we have shown that process $X$ is also stationary. 
\end{proof}
%\end{exmp}
\end{shaded}
\begin{rem}
For a stationary stochastic process, all the existing moments are shift invariant when they exist.  
\end{rem}
\begin{defn}
A \textbf{second order} stochastic process $X$ has finite auto-correlation $R_X(t,t) < \infty$ for all %indices 
$t \in T$. 
\end{defn}
\begin{rem}
This implies $R_X(t_1, t_2) < \infty$ by Cauchy-Schwartz inequality, 
and hence the mean, auto-correlation, and the auto-covariance functions are well defined and finite. 
\end{rem}
\begin{rem}
For a stationary process $X$, we have $X_t = X_0$ and $(X_t,X_s) = (X_{t-s}, X_{0})$ in distribution. 
Therefore, for a second order stationary process $X$, we have 
\meq{3}{
&m_X = \E X_0 , && R_X(t, s) = R_X(t-s,0) = \E X_{t-s}X_0, && C_X(t-s, 0) = R_X(t-s, 0) - m_X^2.
}
\end{rem}
\begin{defn}
A random process $X$ is \textbf{wide sense stationary} if 
\begin{enumerate}
\item $m_X(t) = m_X(t+s)$ for all $s,t \in T$, and 
\item $R_X(t,s) = R_x(t+u, s+u)$ for all $s,t,u \in T$. 
\end{enumerate} 
\end{defn}
\begin{rem}
It follows that a second order stationary stochastic process $X$,  is wide sense stationary. 
A second order wide sense stationary process is not necessarily stationary. 
We can similarly define joint stationarity and joint wide sense stationarity for two stochastic processes $X$ and $Y$.
\end{rem}
\begin{shaded}
\begin{exmp}[Gaussian process] 
Let $X: \Omega \to \R^\R$ be a continuous-time Gaussian process, defined by its finite dimensional distributions. 
In particular, for any finite $S \subset \R$, column vector $x_S \in \R^S$,
%we can define a column vector $x_S$ such that $x_S(s) = x_s$ for $s \in S$, 
mean vector $\mu_S \triangleq \E X_S$, 
and the covariance matrix $C_S \triangleq \E X_SX_S^T$, 
the finite-dimensional density is given by 
\EQ{
f_{X_S}(x_S) = \frac{1}{(2\pi)^{\abs{S}/2}\sqrt{\det(C_S)}}\exp\left(-\frac{1}{2}(x_S-\mu_S)^TC_S^{-1}(x_S-\mu_S)\right).
} 
\end{exmp}
\begin{thm} 
A wide sense stationary Gaussian process is stationary. 
\end{thm}
\proof{
%\begin{rem}
For Gaussian random processes, first and the second moment suffice to get any finite dimensional distribution. 
%\end{rem}
Let $X$ be a wide sense stationary Gaussian process and let $S \subseteq \R$ be finite. 
From the wide sense stationarity of $X$, we have $\E X_S = \mu\bU_S$ and 
\EQ{
\E(X_s-\mu)(X_u-\mu) = C_{s-u},~\text{ for all } s,u \in S. 
}
This means that $C_S = C_{t+S}$, and the result follows.  
}
\end{shaded}

\subsection{Random Walk}
\begin{defn}
Let $X: \Omega \to \sX^\N$ be an \iid random sequence defined on the probability space $(\Omega, \sF, P)$ and the state space $\sX = \R^d$. 
A random sequence $S: \Omega \to \sX^{\Z_+}$ is called a \textbf{random walk} with step-size sequence $X$, 
if $S_0 \triangleq 0$ and $S_n \triangleq \sum_{i=1}^nX_i$ for $n \in \N$. 
\end{defn}
%, then the process $S = (S_n: n \in \N)$ is called a \textbf{random walk}. 
\begin{rem}
We can think of $S_n$ as the random location of a particle after $n$ steps, 
where the particle starts from origin and takes steps of size $X_i$ at the $i$th step. 
From the \iid nature of step-size sequence, 
we observe that $\E S_n = n\E X_1$ and $C_S(n,m) = (n\wedge m)\Var[X_1]$.
\end{rem}
\begin{rem} 
For the process $S: \Omega \to \sX^\N$ it suffices to look at finite dimensional distributions for finite sets $[n] \subseteq \N$ for all $n \in \N$. 
If the \iid step-size sequence $X$ has a common density function, 
then from the transformation of random vectors, 
we can find the finite dimensional density 
\EQ{
f_{S_1, \dots, S_n}(s_1, s_2, \dots, s_n) 
= \frac{1}{\abs{J(s)}}f_{X_1, \dots, X_n}(s_1, s_2-s_1, \dots, s_n-s_{n-1})
= f_{X_1}(s_1)\prod_{i=2}^nf_{X_1}(s_i-s_{i-1}).
}
Recall that $J_{ij}(s) = \frac{\partial s_i}{\partial x_j} = \SetIn{j \le i}$ and hence $J(s)$ is a lower triangular matrix with all non-zero entries being unity, and hence $\abs{J(s)} = 1$.
\end{rem}
\begin{thm} 
The stochastic process $S: \Omega \to \Z_+^{\Z_+}$ has stationary and independent increments.
\end{thm}
\begin{proof}
One needs to show that for $(m_1, \dots, m_n) \subseteq \Z_+$, 
the joint distribution of increments $(S_{m_2}-S_{m_1}, \dots, S_{m_n}-S_{m_{n-1}})$ is independent and distributed identically to $(S_{k+m_2}-S_{k+m_1}, \dots, S_{k+m_n}-S_{k+m_{n-1}})$ for any and $k \in \Z_+$.  
For any $j \in [n-1]$, 
we define event space $\sE_j \triangleq \sigma(X_{m_j+1}, \dots, X_{m_{j+1}})$ and observe that the $j$th increment $S_{m_{j+1}}-S_{m_j}$ is measurable with respect to $j$th event space $\sE_j$. 
Since $X$ is an independent process, it follows that $(\sE_1, \dots, \sE_{n-1})$ are independent event spaces, and hence so are the increments. 
Stationarity of increments follow from the fact that the process $X$ is \iid and $j$th increment $S_{k+m_{j+1}}-S_{k+m_j}$ is sum of $m_{j+1}-m_j$ \iid random variables, and hence an identical distribution for all $k \in \Z_+$. 
\end{proof}
\begin{cor} 
Let $p \in \N$ and for each $i \in [p]$ let $n \in \N^p, k \in \Z_+^p$ such that $n_1 \le \dots \le n_p$ and $k_1 \le \dots \le k_p$. 
%For a finite ordered set $S = (n_1, n_1+n_2, \dots, n_1 + n_2 + \dots + n_p) \subset \N$ and $k_S = (k_1, k_1+k_2, \dots, k_1 + k_2+ \dots + k_p)$, we have 
Then, we can write the joint mass function 
\EQ{
P_{S_{n_1}, \dots, S_{n_p}}(k_1, \dots, k_p) 
= P(\cap_{i \in [k]}\set{S_{n_i} = k_i}) 
= \prod_{i = 1}^{p}P_{S_{n_i-n_{i-1}}}(k_i-k_{i-1}).
}
\end{cor}
\begin{proof}
The result follows from stationary and independent increment property of the random walk $S$. %counting process $S_n$. 
\end{proof}
%%\begin{cor} 
%%Second order moment for the number of success process is given by $R_{N}(n,m) = \E N_n N_m = $
%%\end{cor}
%%\begin{proof}
%%The result follows from stationary and independent increment property of the counting process $N_n$. 
%%\end{proof}


%From previous section, we know following properties of random walks. 
%\begin{thm} 
%For a random walk $(S_n: n \in \N)$ with \emph{iid} step-size sequence $X$, the following are true. 
%\begin{enumerate}[i\_]
%\item The first two moments are $\E S_n = n\E X_i$ and $\Var[S_n] = n\Var[X_i]$. 
%\item Random walk is non-stationary with stationary and independent increments. 
%\item Random walk is homogeneous Markov sequence. 
%\end{enumerate}
%\end{thm}
%For a Bernoulli sequence $X: \Omega \to\set{0,1}^\N$ with $P\set{X_i = 1} = p$, 
%the one dimensional random walk $S$ with step-size sequence $X$ is a positive integer valued random sequence. 
%%with unit step-size. 
\begin{rem} 
For a one-dimensional random walk $S: \Omega\to\Z_+^\N$ with \iid step size sequence $X:\Omega\to\set{0,1}^\N$ such that $P\set{X_1 = 1} = p$, 
the distribution for the random walk at $n$th step $S_n$ is Binomial $(n,p)$. 
That is, 
\EQ{
P\set{S_n = k} = \binom{n}{k}p^k(1-p)^{n-k},\quad k \in \set{0, \dots, n}.
} 
%\begin{enumerate}[i\_]
%%\item The first two moments are $\E S_n = n(p-q)$ and $\Var[S_n] = n(1 - (p-q)^2)$. 
%\item The distribution for the random walk at $n$th step $S_n$ is Binomial $(n,p)$. 
%\item $P\{S_n = k\} = \binom{n}{(n+k)/2}p^{(n+k)/2}q^{(n-k)/2}$ for $n+k$ even, and $0$ otherwise. 
%\end{enumerate}
\end{rem}

%\subsubsection{Independent increments process}
\subsection{L\'evy processes}
A right continuous with left limits stochastic process $X:\Omega\to\R^T$ for index set $T \subseteq \R_+$ with $X_0 = 0$ almost surely, is a \textbf{L\'evy process} if the following conditions hold. 
\begin{enumerate}[(L1)]
%\item $X_{0}=0$ almost surely.
\item The increments are independent. 
For any instants $ 0\le t_{1}<t_{2}<\cdots <t_{n}<\infty$, the random variables $X_{t_{2}}-X_{t_{1}},X_{t_{3}}-X_{t_{2}},\dots ,X_{t_{n}}-X_{t_{n-1}}$ are independent.
\item The increments are stationary. 
For any instants $0\le t_{1}<t_{2}<\cdots <t_{n}<\infty$ and time-difference $s > 0$, 
the random vectors $(X_{t_{2}}-X_{t_{1}},X_{t_{3}}-X_{t_{2}},\dots ,X_{t_{n}}-X_{t_{n-1}})$ and $(X_{s+t_{2}}-X_{s+t_{1}},X_{s+t_{3}}-X_{s+t_{2}},\dots ,X_{s+t_{n}}-X_{s+t_{n-1}})$ are equal in distribution. 
%For any $s<t$,  $X_{t}-X_{s}$, is equal in distribution to $X_{t-s}$.
\item Continuous in probability. For any $\epsilon > 0$ and $t\ge 0$ it holds that $\lim _{h\to 0}P(\set{\abs{X_{t+h}-X_{t}}>\epsilon})=0$. 
\end{enumerate}

\begin{shaded}
\begin{exmp}
Two examples of L\'evy processes are Poisson process and Wiener process. 
The distribution of Poisson process at time $t$ is Poisson with rate $\lambda t$ and the distribution of Wiener process at time $t$ is zero mean Gaussian with variance $t$. 
\end{exmp}

\begin{exmp}
A random walk $S:\Omega\to\sX^{\Z_+}$ with \iid step-size sequence $X:\Omega\to\sX^\N$, is non-stationary with stationary and independent increments. 
To see non-stationarity, we observe that the mean $m_S(n) = n \E X_1$ depends on the step of the random walk. 
We have already seen the increment process of random walks. 
\end{exmp}
%
%\begin{thm} 
%The stochastic process $S: \Omega \to \Z_+^{\Z_+}$ has stationary and independent increments.
%\end{thm}
%\begin{proof}
%We can look at one increment $S_{m+n}-S_m = \sum_{i = 1}^{n}X_{m+i}.$
%This increment is a function of sequence of random variables $(X_{m+1}, \dots, X_{m+n})$ and hence independent of $(X_1, \dots, X_m)$. 
%The random variable $S_m$ depends solely on $(X_1, \dots, X_m)$ and hence the independence follows. 
%Stationarity follows from the fact that the Bernoulli process $X$ is \iid and $S_{m+n}-S_m$ is sum of $n$ \iid Bernoulli random variables, 
%and hence has a Binomial $(n,p)$ distribution identical to that of $S_n$. 
%\end{proof}
\end{shaded}

\subsection{Markov processes}
\begin{defn}
A stochastic process $X$ is \textbf{Markov} if conditioned on the present state, future is independent of the past. 
We denote the history of the process until time $t$ as $\sF_t = \sigma(X_s, s \le t)$. 
That is, for any ordered index set $T$ containing any two indices $u > t$, we have  %such that $t < u$, % \in T$, 
\EQ{
P(\set{X_u \le x_u}\mid \sF_t) = P(\set{X_{u} \le x_u}\mid \sigma(X_t)).
}
The range of the process is called the \textbf{state space}. 
\end{defn}
\begin{rem}
%For a stationary Markov process, we would have 
We next re-write the Markov property more explicitly for the process $X$. 
For all $x, y \in \sX$, finite set $S \subseteq T$ such that $\max S < t < u$, and $H_S  = \cap_{s \in S}\set{X_s \le x_s} \in \sF_t$, we have 
\EQ{
P(\set{X_u \le y}\mid H_S \cap \set{X_t \le x}) = P(\set{X_{u} \le y}\mid \set{X_t \le x}).
}
\end{rem}
\begin{rem}
When the state space $\sX$ is countable, we can write $H_S = \cap_{s \in S}\set{X_s = x_s}$ and the Markov property can be written as 
\EQ{
P(\set{X_u = y}\mid H_S \cap \set{X_t = x}) = P(\set{X_{u} = x_u}\mid \set{X_t = x}).
}
\end{rem}
\begin{rem}
In addition, when the index set is countable, i.e. $T = \Z_+$, then we can take past as $S = \set{0, \dots, n-1}$, 
present as instant $n$, and the future as $n+1$. 
Then, the Markov property can be written as 
\EQ{
P(\set{X_{n+1} = y}\mid H_{n-1} \cap \set{X_n = x}) = P(\set{X_{n+1} = y}\mid \set{X_n = x}),
}
for all $n \in \Z_+, x,y \in \sX$. 
\end{rem}
We will study this process in detail in coming lectures. 
\begin{shaded}
\begin{exmp}
A random walk $S:\Omega\to\sX^{\Z_+}$ with \iid step-size sequence $X:\Omega\to\sX^\N$, is a homogeneous Markov sequence. 
For any $n \in \Z_+$ and $x,y, s_1, \dots, s_{n-1} \in \sX$, we can write the conditional probability 
\EQ{
P(\set{S_{n+1} = y}\mid\set{S_n = x, S_{n-1} = s_{n-1}, \dots, S_1 = s_1}) 
= P(\set{S_{n+1} - S_n = y-x})
=  P(\set{S_{n+1} = y}\mid\set{S_n = x}).
}
\end{exmp}
\begin{lem} 
The stochastic process $S:\Omega\to\Z_+^{\Z_+}$ is homogeneously Markov. 
\end{lem}
\begin{proof} 
Since the process has stationary and independent increments, we have 
\EQ{
P(\set{S_{n+m} = k} \mid\set{S_1 = k_1, S_2  = k_2, \dots, S_n = k_n}) 
= P(\set{S_{n+m}  - S_n = k - k_n}) 
= P(\set{S_{n+m} = k}\mid\set{S_n = k_n}). 
}
\end{proof}
\end{shaded}





\end{document}

